\documentclass[10pt,a4paper]{amsart}
\usepackage[margin=19mm]{geometry}
\usepackage{amsmath,amssymb,mathtools}
\usepackage[scr=rsfs,cal=euler]{mathalfa}
\usepackage{xfrac}
\usepackage[unicode,hidelinks,bookmarks=false]{hyperref}
\newcommand{\ie}{i.e.\ }
\newcommand{\cf}{cf.\ }
\newcommand{\resp}{respectively\ }
\newcommand{\Subsection}{\subsection}
\providecommand{\nobreakdash}{\nobreak\hbox{-}}
\setlength{\emergencystretch}{2em}
\multlinegap0pt
\usepackage{bm}
\usepackage{bbm}
\usepackage{dsfont}
\usepackage{xspace}
\usepackage{cancel}
\usepackage{mathtools}
\usepackage{amsthm}
\makeatletter
\@ifundefined{proposition}{\newtheorem{proposition}{Proposition}}{}
\makeatother
\newcommand{\bd}{\boldsymbol}
\title[Existence and long-time behavior of global strong solutions ]{Existence and long-time behavior of global strong solutions to a nonlinear model of tumor growth}
\author{Jeffrey Kuan, Konstantina Trivisa}
\date{Source: arXiv:2508.19133v1 (CC BY 4.0)}
\begin{document}
\maketitle

\noindent\emph{Source and licence.} Existence and long-time behavior of global strong solutions to a nonlinear model of tumor growth. Version arXiv:2508.19133v1 (CC BY 4.0), distributed under the Creative Commons Attribution 4.0 International licence, as recorded in the arXiv licence metadata for this exact version and shown on the abstract page https://arxiv.org/abs/2508.19133v1. Full rights notice: licenses/arxiv-2508.19133-CC-BY-4.0.txt.\par\smallskip

\noindent\textit{Editorial scope.} Counted unit: the Proposition whose source label is \texttt{regularity} in the article (source0.tex lines 881--887, source label \texttt{regularity}), delivered together with the article's own complete original proof of it (source0.tex lines 889--1049, one \texttt{proof} environment of 161 lines). No mathematics is written, repaired or re-derived: statement and proof are the article's own text line for line. Required context retained uncounted: equations (lines 60--65); CZest (lines 151--160); nonnegativeW (lines 167--177); brinkmann (lines 169--171). Every definition, display and label of those lines is retained unchanged. Omitted from this package and not redistributed: 1--59, 66--150, 161--166, 172--880, 888--888, 1050--1776. Nothing inside a retained excerpt is omitted or altered: every retained body is delivered exactly as the article's lines are written, so no omission receipt of any kind accompanies this package. Citations: the retained excerpts carry no citation command at all, so no bibliography entry of the article is transcribed and no reference is redistributed. Reference adaptation: none was needed and none was made; every reference inside the delivered unit resolves inside it, so no unresolved reference or citation is produced. Printed slips kept literal: no printed word, symbol, hypothesis or number of the counted statement or of its proof was altered. Layout and class: the article's own class is \texttt{article} with packages including biblatex, amsmath, amssymb, amsthm, esint, hyperref, geometry, mathabx, xcolor, graphicx; the wrapper is the lane's pinned \texttt{amsart} editorial wrapper at 10pt on A4, which restates the theorem-like environments the retained text needs and adds the article's own macro definitions that the retained text needs (\texttt{\textbackslash bd}), with extra wrapper packages (bm, bbm, dsfont, xspace, cancel, mathtools). The delivered numbering is the unit's own. Assets: no retained span contains a figure environment, a graphics-inclusion command, a PDF-page inclusion, a table or a TikZ picture; no image file is copied into this package, the package declares no asset dependency and no image file is redistributed with it. Licence: version arXiv:2508.19133v1 of this article is distributed under the Creative Commons Attribution 4.0 International licence, as recorded in the arXiv licence metadata for this exact version and shown on the abstract page https://arxiv.org/abs/2508.19133v1. A full rights notice is delivered at licenses/arxiv-2508.19133-CC-BY-4.0.txt. Characters: every retained excerpt is pure ASCII, so the delivered engine, XeLaTeX with its default Latin Modern text font, reports no missing character.
\begin{equation}\label{equations}
\begin{cases}
    \partial_{t}n - \text{div}(n\nabla W) = \alpha n - \beta n^{\gamma \theta + 1} \\
    -\mu \Delta W + W = an^{\gamma}
\end{cases} \quad \text{ on } \Omega,
\end{equation}

\begin{proposition}\label{CZest}
Suppose that $W$ satisfies the Neumann boundary problem:
\begin{equation*}
-\mu \Delta W + W = f \text{ on } \Omega, \qquad \qquad \nabla W \cdot \bd{n}|_{\partial \Omega} = 0.
\end{equation*}
Then, for all positive integers $m$, 
\begin{equation*}
\|W\|_{W^{m + 2, p}(\Omega)} \le C(m, p) \Big(\|f\|_{W^{m, p}(\Omega)} + \|W\|_{L^{p}(\Omega)}\Big).
\end{equation*}
\end{proposition}

\begin{proposition}\label{nonnegativeW}
    Consider the equation:
    \begin{equation}\label{brinkmann}
        -\mu \Delta W + W = an^{\gamma} \quad \text{ on } \Omega,
    \end{equation}
    with Neumann boundary conditions for a nonnegative $n \in L^{\infty}(0, T; W^{1, p}(\Omega))$ for $p > d$. Then, there exists a unique classical solution $W \in L^{\infty}(0, T; W^{3, p}(\Omega))$ to \eqref{brinkmann} such that:
    \begin{equation*}
        \|W\|_{L^{\infty}(0, T; W^{3, p}(\Omega))} \le C_p\|n\|_{L^{\infty}(0, T; W^{1, p}(\Omega))}^{\gamma},
    \end{equation*}
    for a constant $C_p$ depending on $p$ that is independent of $T$. Also, $W \ge 0$ is nonnegative.
\end{proposition}

\begin{proposition}[Regularity of local-in-time strong solutions]\label{regularity}
Suppose that $(n, W)$ with regularity $n \in C(0, T; W^{1, p}(\Omega))$, $\partial_{t}n \in C(0, T; W^{1, p}(\Omega))$, and $W \in C(0, T; W^{3, p}(\Omega))$ for some positive even integer $p > d$ is a strong solution to \eqref{equations} for smooth initial data $n_{0}(x) > 0$, with $\min_{(x, t) \in [0, T] \times \overline{\Omega}} n(t, x) := n_{min, T} > 0$. Define $\max_{[0, T] \times \overline{\Omega}} n(t, x) := n_{max, T}$. Then, $n(t, \cdot)$ and $\partial_{t}n(t, \cdot)$ are spatially smooth functions for each $t \in [0, T]$ on $\overline{\Omega}$, with the estimates:
\begin{equation*}
\|D^{k}n(t)\|_{L^{p}(\Omega))}^{p} \le e^{C^{*}t}\Big(1 + \|D^{k}n_{0}\|_{L^{p}(\Omega)}^{p}\Big), \qquad \|W(t)\|_{W^{k + 2, p}(\Omega)} \le C^{*}\|n(t)\|_{W^{k, p}(\Omega)},
\end{equation*}
for all positive integers $k$ and for a constant $C^{*}$ depending only on $$C^{*} := C(p, \|n\|_{L^{\infty}(0, T; W^{k - 1, p}(\Omega))}, n_{min, T}, n_{max, T}),$$ which is increasing as a function of $\|n\|_{L^{\infty}(0, T; W^{k - 1, p}(\Omega))}$.
\end{proposition}

\begin{proof}

We show this through a priori estimates. Because we are already assuming that we have a solution on $[0, T]$, we can use bounds for lower derivatives, namely for $n$ in $C(0, T; W^{1, p}(\Omega))$, to help close the \textit{a priori estimates} for the higher derivatives.

\medskip

\noindent \textbf{Estimates on the second derivatives of the density.} We take the second derivative $\partial_{j} \partial_{k}$ of both equations and obtain:
\begin{multline}\label{njk}
\partial_{t} n_{jk} - \sum_{i = 1}^{d} \Big(n_{ijk}W_{i} + n_{ij}W_{ik} + n_{ik}W_{ij} + n_{i}W_{ijk} + n_{jk}W_{ii} + n_{j}W_{iik} + n_{k}W_{iij} + nW_{iijk}\Big) \\
= \alpha n_{jk} - \beta(\gamma \theta + 1) \gamma \theta n^{\gamma \theta - 1} n_{j}n_{k} - \beta(\gamma \theta + 1) n^{\gamma \theta} n_{jk},
\end{multline}
\begin{equation}\label{njk2}
\mu \sum_{i = 1}^{d} W_{iijk} = W_{jk} - a\gamma(\gamma - 1)n^{\gamma - 2} n_{j}n_{k} - a\gamma n^{\gamma - 1} n_{jk}.
\end{equation}
We test the first equation by $p n_{jk}^{p - 1}$ and recall that we so far have estimates for all $1 \le p < \infty$ on $n \in L^{\infty}(0, T; W^{1, p}(\Omega))$ and $W \in L^{\infty}(0, T; W^{3, p}(\Omega))$. After testing \eqref{njk} with $p n_{jk}^{p - 1}$, we obtain the following estimates:
\begin{itemize}
    \item First, we obtain after integration by parts:
    \begin{align*}
    p\Big|\sum_{i = 1}^{d} \int_{0}^{t} \int_{\Omega} \Big(n_{ijk}n_{jk}^{p - 1}W_{i} &+ n_{jk}^{p}W_{ii}\Big)\Big| = (p - 1) \left|\int_{0}^{t} \int_{\Omega} n_{jk}^{p} \Delta W\right| \\
    &\le \mu^{-1}(p - 1) \left|\int_{0}^{t} \int_{\Omega} n_{jk}^{p}W\right| + a\mu^{-1}(p - 1) \left|\int_{0}^{t} n_{jk}^{p}n^{\gamma}\right| \\
    &\le C_{p}\int_{0}^{t} \left(\int_{\Omega} |D^{2}n|^{p}\right) \Big(\|W\|_{L^{\infty}(0, T; H^{2}(\Omega))} + \|n\|_{L^{\infty}(0, T; L^{\infty}(\Omega))}^{\gamma}\Big) \\
    &\le C(p, n_{max, T}) \int_{0}^{t} \int_{\Omega} |D^{2}n|^{p}.
    \end{align*}
    \item Next, we use \eqref{njk2} to estimate
    \begin{align*}
    \Big|p\sum_{i = 1}^{d} &\int_{0}^{t} \int_{\Omega} nn_{jk}^{p - 1}W_{iijk}\Big| \\ &\le C_{p}\left(\left|\int_{0}^{t} \int_{\Omega} nn_{jk}^{p - 1}W_{jk}\right| + \left|\int_{0}^{t} \int_{\Omega} n^{\gamma - 1} n_{j}n_{k} n_{jk}^{p - 1}\right| + \left|\int_{0}^{t} \int_{\Omega} n^{\gamma}n_{jk}^{p}\right|\right) \\
    &\le C_{p}\left(\int_{0}^{t} \int_{\Omega} \Big(n^{p}|W_{jk}|^{p} + n^{(\gamma - 1)p} |\nabla n|^{2p}\Big) + \int_{0}^{t} \int_{\Omega} \Big(1 + n^{\gamma}\Big)|D^{2}n|^{p}\right) \\
    &\le C(p, n_{min, T}, n_{max, T})\left(\int_{0}^{t} \Big(\|W\|^{p}_{W^{2, p}(\Omega)} + \|\nabla n\|_{L^{\infty}(\Omega)}^{p} \|\nabla n\|_{L^{p}(\Omega)}^{p}\Big) + \int_{0}^{t} \int_{\Omega} |D^{2}n|^{p}\right) \\
    &\le C(p, \|n\|_{C(0, T; W^{1, p}(\Omega))}, n_{min, T}, n_{max, T}) \int_{0}^{t} \int_{\Omega}\Big(1 + |D^{2}n|^{p}\Big),
    \end{align*}
    using $p > d$ along with the Sobolev embedding $W^{2, p}(\Omega) \subset W^{1, \infty}(\Omega)$ to estimate $\|\nabla n\|_{L^{\infty}(\Omega)}$ and the assumption that $W \in C(0, T; W^{3, p}(\Omega)) \subset C(0, T; W^{2, \infty}(\Omega))$. Note that by Calder\'{o}n-Zygmund estimates (Proposition \ref{CZest} and \ref{nonnegativeW}), $\|W\|_{C(0, T; W^{2, p}(\Omega))} \le C_{p}\|n\|_{C(0, T; L^{p}(\Omega))}$.
    \item We estimate $\partial_{jk}(n^{\gamma})$ in $L^{p}(\Omega)$ as:
    \begin{align*}
    \|n^{\gamma - 2} n_{j}n_{k} + n^{\gamma - 1} n_{jk}\|_{L^{p}(\Omega)} &\le C(n_{min, T}, n_{max, T}) \Big(\|n_{j}\|_{L^{p}(\Omega)} \|n_{k}\|_{L^{\infty}(\Omega)} + \|n_{jk}\|_{L^{p}(\Omega)}\Big) \\
    &\le C(n_{min, T}, n_{max, T}) \Big(\|n_{j}\|_{L^{p}(\Omega)} \|n\|_{W^{2, p}(\Omega)} + \|n_{jk}\|_{L^{p}(\Omega)}\Big) \\
    &\le C(n_{min, T}, n_{max, T}, \|n\|_{C(0, T; W^{1, p}(\Omega))}) \Big(1 + \|D^{2}n\|_{L^{p}(\Omega)}\Big).
    \end{align*}
    So by Calder\'{o}n-Zygmund estimates (Proposition \ref{CZest} and \ref{nonnegativeW}),
    \begin{equation*}
    \|W\|_{W^{3, \infty}(\Omega)} \le C\|W\|_{W^{4, p}(\Omega)} \le C(p, \|n\|_{C(0, T; W^{1, p}(\Omega))}, n_{min, T}, n_{max, T}) \left(1 + \|D^{2}n\|_{L^{p}(\Omega)}\right).
    \end{equation*}
    Hence, we can estimate:
    \begin{align*}
    \Bigg|\sum_{i = 1}^{d} \int_{0}^{t} \int_{\Omega} &\Big(n_{ij}n_{jk}^{p - 1}W_{ik} + n_{ik}n_{jk}^{p - 1}W_{ij} + n_{i}n_{jk}^{p - 1}W_{ijk} + n_{j}n_{jk}^{p - 1}W_{iik} + n_{k}n_{jk}^{p - 1}W_{iij}\Big)\Bigg| \\
    &\le C_{p}\left(\int_{0}^{t} \int_{\Omega} |\nabla n|^{p} |D^{3}W|^{p} + \int_{0}^{t} \left(\int_{\Omega} |D^{2}n|^{p}\right)\Big(1 +  \|W\|_{L^{\infty}(0, T; W^{2, \infty}(\Omega))}^{p}\Big)\right) \\
    &\le C_{p} \left(\int_{0}^{t} \|\nabla n\|_{L^{\infty}(\Omega)}^{p} \|W\|_{W^{3, p}(\Omega)}^{p} + \int_{0}^{t} \left(\int_{\Omega} |D^{2}n|^{p}\right) \Big(1 + \|W\|_{L^{\infty}(0, T; W^{2, \infty}(\Omega)}^{p}\Big)\right) \\
    &\le C(p, \|n\|_{C(0, T; W^{1, p}(\Omega))}, n_{min, T}, n_{max, T})\int_{0}^{t} \int_{\Omega} \Big(1 + |D^{2}n|^{p}\Big),
    \end{align*}
    since $\|\nabla n\|_{L^{\infty}} \le C_p \|n\|_{W^{2, p}(\Omega)}$ and $W \in C(0, T; W^{3, p}(\Omega)) \subset C(0, T; W^{2, \infty}(\Omega))$ with $$\|W\|_{C(0, T; W^{3, p}(\Omega))} \le C(p, n_{max, T})\|n\|_{C(0, T; W^{1, p}(\Omega))}.$$
    \item Finally, we estimate
    \begin{align*}
    \left|\int_{0}^{t} \int_{\Omega} n^{\gamma \theta - 1} n_{j}n_{k} n_{jk}^{p - 1}\right| &\le C(n_{min, T}, n_{max, T}) \int_{0}^{t} \|n_{k}\|_{L^{\infty}(\Omega)} \|n_{j}\|_{L^{p}(\Omega)} \|n_{jk}\|_{L^{p}(\Omega)}^{p - 1} dt \\
    &\le C(\|n\|_{C(0, T; W^{1, p}(\Omega))}, n_{min, T}, n_{max, T}) \int_{0}^{t} \int_{\Omega} \Big(1 + |D^{2}n|^{p}\Big),
    \end{align*}
    by the estimate $\|n_{k}\|_{L^{\infty}(\Omega)} \le C\|n\|_{W^{2, p}(\Omega)}$. We also have the estimate:
    \begin{equation*}
    \beta p (\gamma \theta + 1) \left|\int_{0}^{t} \int_{\Omega} n^{\gamma \theta} n_{jk}^{p}\right| \le C(n_{min, T}, n_{max, T}) \int_{0}^{t} \int_{\Omega} |D^{2}n|^{p}.
    \end{equation*}
\end{itemize}

\noindent \textit{Conclusion.} Combining all of the above estimates, we obtain:
\begin{equation*}
\int_{\Omega} |D^{2}n(t)|^{p} \le \int_{\Omega} |D^{2}n_{0}|^{p} + C(p, \|n\|_{C(0, T; W^{1, p}(\Omega))}, n_{min, T}, n_{max, T}) \int_{0}^{t} \int_{\Omega} \Big(1 + |D^{2}n(s)|^{p}\Big) ds,
\end{equation*}
so by Gronwall's inequality and Calder\'{o}n-Zygmund estimates, we obtain:
\begin{equation*}
\|D^{2}n(t)\|_{L^{p}(\Omega)}^{p} \le e^{C^{*}t}\Big(1 + \|D^{2}n_{0}\|_{L^{p}(\Omega)}^{p}\Big), \qquad \|W(t)\|_{W^{4, p}(\Omega)} \le C^{*}\|n\|_{W^{2, p}(\Omega)},
\end{equation*} 
for some constant $C^{*}$ depending on the parameters $C^{*} := C(p, \|n\|_{C(0, T; W^{1, p}(\Omega))}, n_{min, T}, n_{max, T})$.

\medskip

\noindent \textbf{Estimates on all higher derivatives.} We obtain estimates on higher order spatial derivatives by using an inductive argument, with an inductive assumption that for some positive integer $k$, we have established the following bound:
\begin{equation}\label{inductivek}
\|n\|_{L^{\infty}(0, T; W^{k, p}(\Omega))} \le C, \quad \text{ and hence, by Calder\'{o}n-Zygmund}, \quad \|W\|_{L^{\infty}(0, T; W^{k + 2, p}(\Omega))} \le C.
\end{equation}
We then want to show that under this inductive assumption, we can obtain the estimates:
\begin{equation}\label{inductivestep}
\|D^{k + 1}n\|_{L^{p}(\Omega)}^{p} \le e^{C^{*}t}\Big(1 + \|D^{k + 1}n_{0}\|_{L^{p}(\Omega)}^{p}\Big), \qquad \|W(t)\|_{W^{k + 3, p}(\Omega)} \le C^{*}\|n(t)\|_{W^{k + 1, p}(\Omega)},
\end{equation}
for some constant $C^{*}$, where for simplicity of notation, we will use this shorthand to indicate the following explicit dependence of constants on parameters for the remainder of the proof:
\begin{equation}\label{Cstardepend}
C^{*} = C(p, \|n\|_{L^{\infty}(0, T; W^{k, p}(\Omega))}, n_{min, T}, n_{max, T}).
\end{equation}

To do this, we consider a general multi-index $\partial_{\alpha}$ with $|\alpha| = k + 1$ and we take the derivative with respect to this multi-index of both equations. We obtain:
\begin{equation}\label{alphaderivative}
\partial_{t}\partial_{\alpha} n - \sum_{i = 1}^{d} \Big((\partial_{\alpha} \partial_{i} n) \partial_{i}W + (\partial_{\alpha} n) \partial_{i}^{2} W + n\partial_{\alpha} \partial^{2}_{i}W\Big) + R_{\alpha} = \alpha \partial_{\alpha}n - \beta(\gamma \theta + 1)n^{\gamma \theta} \partial_{\alpha} n + T_{\alpha},
\end{equation}
where:
\begin{itemize}
    \item $R_{\alpha}$ is comprised of terms $(\partial_{\beta} n)(\partial_{\gamma} W)$ for multi-indices with $1 \le |\beta| \le k + 1 = |\alpha|$ and $|\gamma| = k + 3 - |\beta|$.
    \item $T_{\alpha}$ is comprised of terms $n^{\gamma \theta - m} (\partial_{\alpha_{1}}n) (\partial_{\alpha_{2}}n) \cdot \cdot \cdot (\partial_{\alpha_{m+1}}n)$, for a positive integer $1 \le m \le k$ and multi-indices $1 \le |\alpha_{1}| \le |\alpha_{2}| \le ... \le |\alpha_{m+1}| \le k - m + 1$ with $\displaystyle \sum_{i = 1}^{m+1} |\alpha_{i}| = k + 1$.
\end{itemize}
We also have the equation for the potential:
\begin{equation}\label{Winductive}
\mu \sum_{i = 1}^{d} \partial_{\alpha} \partial^{2}_{i}W = \partial_{\alpha} W - a \partial_{\alpha} (n^{\gamma}).
\end{equation}
We test \eqref{alphaderivative} with $p(\partial_{\alpha}n)^{p - 1}$ and integrate over space, and estimate the resulting terms.
\begin{itemize}
\item By integrating by parts and using the Neumann boundary condition on $W$:
\begin{equation*}
p\sum_{i = 1}^{d} \int_{0}^{t} \int_{\Omega} \Big((\partial_{\alpha} \partial_{i}n) \partial_{i}W + (\partial_{\alpha} n)\partial^{2}_{i}W\Big) (\partial_{\alpha} n)^{p - 1} = (p - 1)\int_{0}^{t} \int_{\Omega} (\partial_{\alpha} n)^{p} \Delta W. 
\end{equation*}
Since $k + 2 \ge 4$ in the inductive assumption, we can estimate that $\|\Delta W\|_{L^{\infty}(\Omega)} \le \|W\|_{W^{4, p}(\Omega)}$, and hence by the boundedness of $W$ in $L^{\infty}(0, T; W^{k + 2, p}(\Omega))$:
\begin{equation*}
\left|p\sum_{i = 1}^{d} \int_{0}^{t} \int_{\Omega} \Big((\partial_{\alpha} \partial_{i}n) \partial_{i}W + (\partial_{\alpha}n)\partial^{2}_{i}W\Big) (\partial_{\alpha}n)^{p - 1}\right| \le C^{*}\int_{0}^{t} \int_{\Omega} (\partial_{\alpha}n)^{p},
\end{equation*}
where we recall the shorthand $C^{*}$ from \eqref{Cstardepend}.
\item Next, using the equation for the potential \eqref{Winductive}, we estimate that
\begin{equation}\label{midest}
\left|p\sum_{i = 1}^{d} \int_{0}^{t} \int_{\Omega} n(\partial_{\alpha}n)^{p - 1} \partial_{\alpha}\partial^{2}_{i}W\right| \le C_{p}\left(\left|\int_{0}^{t} \int_{\Omega} n(\partial_{\alpha}n)^{p - 1} \partial_{\alpha} W\right| + \left|\int_{0}^{t} \int_{\Omega} n(\partial_{\alpha}n)^{p - 1} \partial_{\alpha}(n^{\gamma})\right|\right).
\end{equation}
For the first term on the right-hand side of \eqref{midest}, using the inductive assumption on $W$ in \eqref{inductivek}:
\begin{multline*}
\left|\int_{0}^{t} \int_{\Omega} n(\partial_{\alpha} n)^{p - 1} \partial_{\alpha} W\right| \le C_{p}\left(\int_{0}^{t} \int_{\Omega} n^{p} |\partial_{\alpha} W|^{p} + \int_{0}^{t} \int_{\Omega} |\partial_{\alpha}n|^{p}\right) \\
\le C_{p}\left(\int_{0}^{t}\|n\|_{L^{\infty}(0, T; L^{\infty}(\Omega))}^{p} \|W\|_{L^{\infty}(0, T; W^{k + 1, p}(\Omega))}^{p} + \int_{0}^{t} \int_{\Omega} |\partial_{\alpha} n|^{p}\right) \le C^{*}\int_{0}^{t} \int_{\Omega} \Big(1 + |\partial_{\alpha} n|^{p}\Big),
\end{multline*}
where we use the shorthand \eqref{Cstardepend} to indicate the dependence of constants on parameters. For the second term on the right-hand side of \eqref{midest}, we note that because $n \in L^{\infty}(0, T; W^{k, p}(\Omega))$ for some $k \ge 2$ (from the inductive assumption \eqref{inductivek}) and for $p > d$, we have that there are $L^{\infty}(\Omega)$ bounds on the first derivative of $n$ and all higher order derivatives up to order $k - 1$. So therefore,
\begin{equation*}
\left|\int_{0}^{t} \int_{\Omega} n(\partial_{\alpha}n)^{p - 1} \partial_{\alpha}(n^{\gamma})\right| \le \int_{0}^{t} \int_{\Omega} \gamma n^{\gamma} |\partial_{\alpha}n|^{p} + \text{l.o.t.} \le C^{*}\int_{0}^{t} \int_{\Omega} \Big(1 + |D^{k + 1}n|^{p}\Big),
\end{equation*}
where the lower order terms are terms of the form $\displaystyle \int_{0}^{t} \int_{\Omega} n^{\gamma - m} |\partial_{\alpha}n|^{p - 1} |\partial_{\alpha_{1}}n \partial_{\alpha_{2}}n \cdot \cdot \cdot \partial_{\alpha_{m+1}}n|$ where $1 \le m \le k$ and $1 \le |\alpha_{1}| \le |\alpha_{2}| \le ... \le |\alpha_{m+1}| \le k - m + 1$ with $|\alpha_{1}| + |\alpha_{2}| + ... + |\alpha_{m+1}| = k + 1$. The bound on the first term $\displaystyle \int_{0}^{t} \int_{\Omega} \gamma n^{\gamma} |\partial_{\alpha}n|^{p}$ follows from the previous bounds on $\|n\|_{L^{\infty}(0, T; L^{\infty}(\Omega))}$. We can then estimate the lower order terms by estimating that 
\begin{align}\label{lotest}
\int_{0}^{t} \int_{\Omega} n^{\gamma - m} |\partial_{\alpha}n|^{p - 1} |\partial_{\alpha_{1}}n \cdot \cdot \cdot \partial_{\alpha_{m+1}}n| &\le C_{p}\left(\int_{0}^{t} \int_{\Omega} |\partial_{\alpha}n|^{p} + \int_{0}^{t} \int_{\Omega} |n^{p(\gamma - m)}| \cdot  |\partial_{\alpha_{1}}n \cdot \cdot \cdot \partial_{\alpha_{m+1}}n|^{p}\right) \nonumber \\
&\le C^{*}\int_{0}^{t} \int_{\Omega} \Big(1 + |D^{k+1}n|^{p}\Big),
\end{align}
where the estimate subsuming the second integral into $C^{*}$, see the definition of $C^{*}$ in \eqref{Cstardepend}, follows from the inductive assumption that $n \in L^{\infty}(0, T; W^{k, p}(\Omega))$ for some $k \ge 2$ and $p > d$. We also observe that the second integral can only contain at most one $|\alpha_{i}|$ equal to $k$ (which can be estimated in $L^{p}(\Omega)$ by the inductive assumption \eqref{inductivek}), in which case the rest of the $|\alpha_{i}|$ are strictly less than $k$ and hence the $\alpha_{i}$-derivatives can be bounded in $L^{\infty}$ using the inductive bound on $n \in L^{\infty}(0, T; W^{k, p}(\Omega))$.

\item Next, we estimate the contribution of $\displaystyle \int_{0}^{t} \int_{\Omega} R_{\alpha} (\partial_{\alpha}n)^{p - 1}$, where we recall that $R_{\alpha}$ is made up of terms of the form $(\partial_{\beta} n)(\partial_{\gamma} W)$ where $1 \le |\beta| \le k + 1$ and $|\gamma| = k + 3 - |\beta|$. If $|\beta| = k + 1$ so that $|\gamma| = 2$, we can estimate:
\begin{equation*}
\left|\int_{0}^{t} \int_{\Omega} \partial_{\beta} n (\partial_{\alpha}n)^{p - 1} \partial_{\gamma} W\right| \le \|W\|_{L^{\infty}(0, T; W^{2, \infty}(\Omega))} \int_{0}^{t} \int_{\Omega} |D^{k + 1}n|^{p} \le C^{*} \int_{0}^{t} \int_{\Omega} |D^{k + 1}n|^{p},
\end{equation*}
by Sobolev embedding and the inductive assumption $W \in L^{\infty}(0, T; W^{k + 2, p}(\Omega))$ for some $k \ge 2$ and $p > d$ in \eqref{inductivek}. In the case where $|\beta| = 1$, then $|\gamma| = k + 2$, and we can use the inductive bound on $\|n\|_{L^{\infty}(0, T; W^{k, p}(\Omega))}$ and $\|W\|_{L^{\infty}(0, T; W^{k + 2, p}(\Omega))}$ for $k \ge 2$, in \eqref{inductivek} to estimate:
\begin{align*}
\left|\int_{0}^{t} \int_{\Omega} \partial_{\beta} n (\partial_{\alpha} n)^{p - 1}\partial_{\gamma} W\right| &\le C_{p}\left(\int_{0}^{t} \int_{\Omega} |\nabla n|^{p} |D^{k + 2}W|^{p} + \int_{0}^{t} \int_{\Omega} |D^{k + 1}n|^{p}\right) \\
&\le C^{*} \int_{0}^{t} \int_{\Omega} \Big(1 + |D^{k + 1}n|^{p}\Big). 
\end{align*}
Finally, if $2 \le |\beta| \le k$ and hence $3 \le |\gamma| \le k + 1$ with $|\gamma| = k + 3 - |\beta|$, we estimate:
\begin{align*}
\left|\int_{0}^{t} \int_{\Omega} (\partial_{\beta} n)(\partial_{\alpha}n)^{p - 1} \partial_{\gamma} W\right| &\le C_p\left(\int_{0}^{t} \int_{\Omega} |\partial_{\beta}n|^{p} |\partial_{\gamma} W|^{p} + \int_{0}^{t} \int_{\Omega} |D^{k + 1}n|^{p}\right) \\
&\le C^{*}\int_{0}^{t} \int_{\Omega} \Big(1 + |D^{k + 1}n|^{p}\Big),
\end{align*}
since by the inductive assumption \eqref{inductivek}, $n \in L^{\infty}(0, T; W^{k, p}(\Omega))$ and $W \in L^{\infty}(0, T; W^{k + 2, p}(\Omega))$ for all $p > d$ and for some $k \ge 2$. So we conclude that the total contribution is:
\begin{equation*}
\left|\int_{0}^{t} \int_{\Omega} R_{\alpha} (\partial_{\alpha} n)^{p - 1}\right| \le C^{*} \int_{0}^{t} \int_{\Omega} \Big(1 + |D^{k + 1}n|^{p}\Big).
\end{equation*}
\item Next, we estimate $\displaystyle \beta (\gamma \theta + 1) p \left|\int_{0}^{t}  \int_{\Omega} n^{\gamma \theta} (\partial_{\alpha}n)^{p}\right| \le C^{*} \int_{0}^{t} \int_{\Omega} |D^{k + 1}n|^{p}$.
\item Finally, we estimate $\displaystyle \left|\int_{0}^{t} \int_{\Omega} T_{\alpha} (\partial_{\alpha}n)^{p - 1}\right|$, where we recall that $T_{\alpha}$ is comprised of terms of the form $n^{\gamma \theta - m}(\partial_{\alpha_{1}}n)(\partial_{\alpha_{2}}n) \cdot \cdot \cdot (\partial_{\alpha_{m + 1}}n)$ for $1 \le m \le k$ and $1 \le |\alpha_{1}| \le |\alpha_{2}| \le ... \le |\alpha_{m + 1}| \le k - m + 1$ with $\displaystyle \sum_{i = 1}^{m + 1} |\alpha_{i}| = k + 1$. As in \eqref{lotest}, we estimate for $C^{*}$ in \eqref{Cstardepend}:
\begin{equation*}
\left|\int_{0}^{t} \int_{\Omega} T_{\alpha}(\partial_{\alpha}n)^{p - 1}\right| \le C^{*}\int_{0}^{t} \int_{\Omega} \Big(1 + |D^{k + 1}n|^{p}\Big).
\end{equation*}
\end{itemize}

So we obtain the final estimate $\displaystyle 
\int_{\Omega} |D^{k + 1}n(t)|^{p} \le \int_{\Omega} |D^{k + 1}n_{0}|^{p} + C^{*} \int_{0}^{t} \int_{\Omega} \Big(1 + |D^{k + 1}n(s)|^{p}\Big)$, and hence conclude the inductive step \eqref{inductivestep} for $\|n\|_{L^{\infty}(0, T; W^{k + 1, p}(\Omega))}$ by Gronwall's inequality and for $\|W\|_{L^{\infty}(0, T; W^{k + 3, p}(\Omega))}$ by Calder\'{o}n-Zygmund estimates, obtaining that for some constant $C^{*}$ depending on parameters as in \eqref{Cstardepend}:
\begin{equation*}
\|D^{k + 1}n(t)\|_{L^{p}(\Omega)}^{p} \le e^{C^{*}t}\Big(1 + \|D^{k + 1} n_{0}\|_{L^{p}(\Omega)}^{p}\Big), \qquad \|W(t)\|_{W^{k + 3, p}(\Omega)} \le C^{*}\|n(t)\|_{W^{k + 1, p}(\Omega)}.
\end{equation*}

\end{proof}

\end{document}
